\let\negmedspace\undefined
\let\negthickspace\undefined
\documentclass[journal]{IEEEtran}
\usepackage[a5paper, margin=10mm, onecolumn]{geometry}
%\usepackage{lmodern} % Ensure lmodern is loaded for pdflatex
\usepackage{tfrupee} % Include tfrupee package

\setlength{\headheight}{1cm} % Set the height of the header box
\setlength{\headsep}{0mm}     % Set the distance between the header box and the top of the text

\usepackage{gvv-book}
\usepackage{gvv}
\usepackage{cite}
\usepackage{amsmath,amssymb,amsfonts,amsthm}
\usepackage{algorithmic}
\usepackage{graphicx}
\usepackage{textcomp}
\usepackage{xcolor}
\usepackage{txfonts}
\usepackage{listings}
\usepackage{enumitem}
\usepackage{mathtools}
\usepackage{gensymb}
\usepackage{comment}
\usepackage[breaklinks=true]{hyperref}
\usepackage{tkz-euclide} 
\usepackage{listings}
% \usepackage{gvv}                                        
\def\inputGnumericTable{}                                 
\usepackage[latin1]{inputenc}                                
\usepackage{color}                                            
\usepackage{array}                                            
\usepackage{longtable}                                       
\usepackage{calc}                                             
\usepackage{multirow}                                         
\usepackage{hhline}                                           
\usepackage{ifthen}                                           
\usepackage{lscape}
\begin{document}

\bibliographystyle{IEEEtran}
\vspace{3cm}

\title{1.8.28}
\author{EE25BTECH11011 - Navya}
% \maketitle
% \newpage
% \bigskip
{\let\newpage\relax\maketitle}

\renewcommand{\thefigure}{\theenumi}
\renewcommand{\thetable}{\theenumi}
\setlength{\intextsep}{10pt} % Space between text and floats

\textbf{Question}:\\
Find the equation of the set of the points $P$ such that its distances 
from the points $A(3,4,-5)$ and $B(-2,1,4)$ are equal.\\

\textbf{Solution}:\\
\begin{align}
\text{Let }A=(3,4,-5),\quad B=(-2,1,4),\quad P=(x,y,z).
\end{align}

Using vector notation $\|\cdot\|$ for Euclidean norm,

\begin{align}
\|P-A\|^2+\|P-B\|^2=2k^2.
\end{align}

Write $M=\dfrac{A+B}{2}$ (the midpoint of $AB$).\\

Using the identity
\begin{align}
\|P-A\|^2+\|P-B\|^2=2\|P-M\|^2+2\|A-M\|^2,
\end{align}

the condition becomes
\begin{align}
2\|P-M\|^2+2\|A-M\|^2=2k^2\\
\|P-M\|^2=k^2-\|A-M\|^2
\end{align}
Now compute $M$ and $\|A-M\|^2$:

\begin{align}
M=\frac{A+B}{2}=\left(\frac{3+(-2)}{2},\frac{4+1}{2},\frac{-5+4}{2}\right)
=\left(\frac{1}{2},\frac{5}{2},-\frac{1}{2}\right),
\end{align}

\begin{align}
AB=B-A=(-5,-3,9),\\
\|AB\|^2=25+9+81=115,
\end{align}

\begin{align}
\|A-M\|^2=\frac{\|AB\|^2}{4}=\frac{115}{4}.
\end{align}

Hence the locus of $P$ is the sphere with centre $M$ and radius
$\sqrt{k^2-\tfrac{115}{4}}$. Its equation is
\begin{align}
\left(x-\tfrac12\right)^2+\left(y-\tfrac52\right)^2+\left(z+\tfrac12\right)^2
= k^2-\frac{115}{4},
\end{align}
which exists (real radius) iff $k^2\ge\frac{115}{4}$
\begin{figure}[h!]
    \centering
    \includegraphics[width=0.6\textwidth]{figs/fig3.png}
    \caption{locus of point p}
    \label{fig:sphere-locus}
\end{figure}
\end{document}